\clearpage
\chapter*{Đính chính cho Bài trình bày số 190}
\addcontentsline{toc}{chapter}{Đính chính cho Bài trình bày số 190}
\setcounter{footnote}{0}
\sid{e190}
\markboth{Đính chính 190}{Đính chính 190}
\src{300}

\noindent\textbf{3. Kỹ thuật hạ xuống và các định lý tồn tại trong hình học
đại số [TDTE I]. I: Đại cương. Hạ xuống theo các cấu xạ phẳng trung thành.}

\begin{center}
[Hội thảo Bourbaki, tập 12, 1959/60, số 190, 29 tr.]
\end{center}

\noindent\textbf{Trang 1, dòng 10.} --- Nay xem ra nói rằng kỹ thuật hạ xuống
«là nền tảng của phần lớn các định lý tồn tại trong hình học đại số» là quá
mức. Điều này đúng phần lớn đối với các kỹ thuật không xạ ảnh được trình bày
trong hai bài đầu của loạt TDTE I đến VI, nhưng không đúng đối với các kỹ
thuật xạ ảnh (TDTE IV, V, VI).

\noindent\textbf{Trang 5, dòng 16.} --- Không cần giả sử rằng $\alpha$ là một
cấu xạ $\mathcal F$-hạ xuống.

\noindent\textbf{Trang 20, nhận xét.} --- Một cấu xạ tựa hữu hạn, étale và
toàn ánh $S'\longrightarrow S$, hoặc một cấu xạ
$\operatorname{Spec}(\widehat A)\longrightarrow\operatorname{Spec}(A)$, không
nhất thiết luôn là một cấu xạ hạ xuống nghiêm ngặt, ngay cả khi $A$ là vành
địa phương của một đường cong đại số trên một trường đóng đại số $k$ và
$S=\operatorname{Spec}(A)$. Do đó có thể tìm được một cấu xạ riêng và trơn
$f:X\longrightarrow S$, làm cho $X$ thành
một phân thớ chính trên $S$ dưới tác động của một đường cong elliptic $E$,
sao cho $f':X'\longrightarrow S'$ là xạ ảnh nhưng $f$ không xạ ảnh. Vì vậy,
đây đồng thời là một ví dụ về phân thớ chính thuần nhất không đẳng tầm dưới
tác động của một lược đồ abel.

\noindent\textbf{Trang 26, dòng 9.} --- Đọc «Chow--Iang» thay cho «Igusa--Lang».

Về các chi tiết khác nhau của lý thuyết hạ xuống, xem SGA VI, VII và VIII.

\par
